\section{Harmonic coordinates and the metric on K3 bundle moduli}
\label{sec:harmonic}

Let $X$ be a projective K3 surface and $H$ an integral ample class.
Let $F_1,\ldots,F_l$ be pairwise nonisomorphic locally free
$H$-slope-stable sheaves with the same reduced $H$-Hilbert polynomial.
Fix positive multiplicities $m_i$ and put
\[
 E_0=\bigoplus_i F_i\otimes W_i,\qquad W_i=\CC^{m_i}.
\]
For the central analytic background choose a real K\"ahler class
$\omega_*$ for which the same $F_i$ are slope-stable of a common slope,
and use its Ricci-flat K\"ahler metric. One may take $\omega_*=H$.
These are separate hypotheses: $H$ supplies algebraic formality,
whereas $\omega_*$ supplies the Hermitian--Einstein background.
Choose its direct-sum Hermitian--Einstein metric $h_0$ on $E_0$.
Fix the determinant, and let $\omega_t$ be a smooth finite-dimensional
family of K\"ahler forms with $\omega_0=\omega_*$.

Give $V=\Ext^1(E_0,E_0)$ its central harmonic Hermitian norm.
Serre duality identifies it unitarily with a doubled-quiver
representation for $K=\prod_iU(W_i)$; self-extensions give paired
loops, as proved below. Let $\mu_{\CC},\mu$ be its complex and real
moments with the normalization of \eqref{eq:moment-normalization},
and put $Z=\mu_{\CC}^{-1}(0)$.
For $r_i=\rank F_i$ and $N=\sum_i r_im_i$, define
\begin{equation}\label{eq:geometric-parameter}
 \zeta_i(t)=2\pi\bigl(r_i\mu_{\omega_t}(E_0)
                  -\deg_{\omega_t}F_i\bigr),
 \qquad \sum_i m_i\zeta_i(t)=0.
\end{equation}
Degrees mean $\int_Xc_1(F_i)\wedge\omega_t$ and slope is degree
divided by rank. In particular \eqref{eq:geometric-parameter}
has no volume denominator.
This sign agrees with the incoming-minus-outgoing convention:
a subbundle of smaller slope has positive total $\zeta$-weight.

At a balanced quiver point, decompose the representation into
pairwise nonisomorphic stable summands. Its \emph{representation type}
is the unordered list of their dimension vectors and multiplicities.
Stability here means strictly positive $\zeta$-weight on every
nonzero proper subrepresentation; Proposition~\ref{prop:real-polystability}
justifies this decomposition. The connected strata of a fixed type
are the orbit-type strata for the compact group action. We use
these strata, including singular types, throughout the metric statements.

\begin{theorem}\label{thm:harmonic-comparison}
Under these hypotheses there are common nested neighborhoods in
which the finite quotients
\[
 M_{\zeta(t)}=(Z\cap\mu^{-1}(\zeta(t)))/K
\]
identify with open neighborhoods in the actual fixed-determinant
Hermitian--Yang--Mills bundle moduli. Every sufficiently small actual
solution for sufficiently small $t$ is covered. The maps $f_t$ are
homeomorphisms, smooth on representation-type strata, and preserve
their complex structures. On genuine stratum tangents at a finite
representative $x$, with $r=\|x\|$,
\begin{equation}\label{eq:harmonic-comparison}
 (1-C(|t|+r))g_{\zeta(t)}
 \le f_t^*g_t^{L^2}
 \le(1+C(|t|+r))g_{\zeta(t)}.
\end{equation}
The constants are independent of stabilizer type and of ratios of
positive eigenvalues of the finite orbit operator.

If a central unitary automorphism acts as $-I$ on all of $V$, the
spatial error in \eqref{eq:harmonic-comparison} improves to $Cr^2$.
At the central value the construction identifies the full reduced
analytic quotient germ, including its singular locus.
\end{theorem}

Here a neighborhood of actual solutions means solutions represented
by small deformations of the fixed smooth bundle and central
Dolbeault operator, uniformly in $t$. The central holomorphic bundle
itself need not remain polystable when $t\ne0$.
The proof combines a fixed elliptic inverse with
Theorem~\ref{thm:uniform-balancing}. The elliptic inverse removes the
whole parallel algebra; no potentially small finite symmetry is
included in that inverse. The complementary minimization method
of \cite{Wue26b,Wue26c,Wue26d} then yields the metric estimate.

\subsection{Unitary quiver coordinates}
Write $d=\bar\partial_0$ on $\End E_0$ and identify $V$ with its
harmonic degree-one forms. Fix a parallel holomorphic two-form
$\Omega_X$ with normalization compatible with the harmonic norm.
In unitary coframes and factor frames define
\[
 J_{ij}(A_1\bar\theta^1+A_2\bar\theta^2)
       =-A_2^*\bar\theta^1+A_1^*\bar\theta^2.
\]
This is frame-independent, antiunitary, and maps the $(i,j)$ factor
space to the $(j,i)$ space, with $J_{ji}J_{ij}=-I$.
Choose the orientation of the Serre pairing
$Q(\alpha,\gamma)=\int_X\Omega_X\wedge\tr(\alpha\wedge\gamma)$
so that
\[
 Q(\alpha,J\beta)=\langle\alpha,\beta\rangle_{L^2}.
\]
Pointwise this is the identity
$\tr(A_1B_1^*+A_2B_2^*)$.
If $\beta$ is harmonic, its orthogonality to exact forms and Stokes'
formula show $dJ\beta=0$. Apply the opposite pairing, $J^2=-I$ and
Stokes' formula to show that $J\beta$ is orthogonal to exact forms
as well. Thus $J$ preserves harmonicity; flat factor connections
are not assumed.

Paired orthonormal bases identify off-diagonal extension spaces
with opposite arrows. Each diagonal $\Ext^1(F_i,F_i)$ is
quaternionic; choose pairs $\alpha,J\alpha$, each giving loop
coefficients $A,B\in\End W_i$. The quadratic Yoneda contribution
of this pair is $[A,B]$. In the Serre-dual target paired with the
parallel algebra by ordinary multiplicity trace, the total quadratic
obstruction is exactly $\mu_{\CC}$.
If $\eta_X$ generates $H^{0,2}(\mathcal O_X)$ with
$\int_X\Omega_X\wedge\eta_X=1$, its literal harmonic block
representative is instead
\begin{equation}\label{eq:rank-dual}
 \bigoplus_i r_i^{-1} I_{F_i}\otimes\mu_{\CC,i}\,\eta_X.
\end{equation}
The factors $r_i^{-1}$ compensate for tracing $I_{F_i}$.
The ordinary finite dual and these harmonic block coefficients
must not be interchanged.

\subsection{Convergent quadratic integrability coordinates}
The harmonic Kuranishi embedding on a full ball is the convergent
solution of
\begin{equation}\label{eq:kuranishi-embedding}
 \Phi(b)=b-d^*G_2(\Phi(b)\wedge\Phi(b)),\qquad D\Phi(0)=\iota.
\end{equation}
Here $G_2$ is the Green operator for the degree-two Dolbeault
Laplacian at the central bundle. The elliptic contraction construction is holomorphic, and its
Green operators commute with the parallel multiplicity action.
Its integrability locus has a reduced analytic germ $S\subset V$.
Reduced analytic semiuniversality is provided by
\cite[\S5 and Remark~4]{BS21}.

The bundle $E_0$ is $H$-Gieseker polystable. Budur--Zhang
\cite[Theorem~1.2]{BZ19} gives formality of its derived
endomorphism algebra. The equivariant analytic result of
Arbarello--Sacc\`a \cite[Proposition~4.4]{AS18}, with the quiver
identification in their Proposition~6.1, identifies the reduced
Quot slice with the quadratic germ $Z$, tangent to the identity.
This is a fixed-bundle germ result. We do not invoke their
dimension-one sheaf wall-crossing theorem for vector bundles.

We record the analytic comparison with the harmonic family,
because formality alone would not supply controlled harmonic
coordinates. Twist the harmonic family by a sufficiently large
multiple of $H$. Analytic base change makes its spaces of sections
a vector bundle with surjective evaluation map
\cite[Corollaries~1.4--1.5 and Theorem~3.1]{Gro61a}.
Average a holomorphic frame over $K$; it remains a frame near the
origin and makes the evaluation equivariant.
The analytic Quot construction \cite[Theorem~1.1]{Gro61b} gives a
classifying map. An equivariant complement to the stabilizer in
the ambient general linear algebra, followed by Luna's slice
\cite{Lun73}, gives a $K$-equivariant germ map
$h:S\to Z_{\rm Quot}$ with identity Kodaira--Spencer derivative.

The reduced tangent space of $Z$ is all of $V$: every subspace
supported on a single arrow, including a single loop with its
paired loop zero, lies in $Z$, and these subspaces span $V$.
First apply reduced semiuniversality to the Quot family to obtain
a reverse map $g:Z_{\rm Quot}\to S$. Its Kodaira--Spencer derivative,
followed by harmonic inclusion, is the identity on $V$.
Thus the reduced tangent of $S$ is also $V$.
This argument does not assert that an arbitrary matrix-valued
harmonic self-extension squares to zero as a differential form.

Both compositions have invertible tangent maps and are analytic
germ isomorphisms. For completeness, lift a self-map with invertible
tangent map to an ambient local biholomorphism in a minimal embedding.
Its pullback on the germ's local ring is surjective.
A surjective endomorphism of a Noetherian ring is injective:
the ascending kernels of its powers stabilize, and surjectivity
then forces its first kernel to vanish. This proves the assertion,
and hence $h$ is an isomorphism.

Extend the resulting coordinate functions to an ambient holomorphic
map and average over $K$. We obtain an equivariant local
biholomorphism $F$ satisfying
\begin{equation}\label{eq:ambient-coordinates}
 \begin{aligned}
 F(x)&=x+O(\|x\|^2),& DF(x)&=I+O(\|x\|),\\
 a(x)&=\Phi(F(x)),& D_x&=d+[a(x),\cdot],\qquad J_x=Da(x).
 \end{aligned}
\end{equation}
Its integrable locus is exactly $Z$ on suitably small representatives.
Compact equivariance complexifies locally, and for every parallel
complex endomorphism $\xi$ gives the exact identity
\begin{equation}\label{eq:exact-gauge-derivative}
 D_x\xi=-J_xC_x\xi.
\end{equation}
For a genuine tangent $v$ to an integrable stratum,
$D_x(J_xv)=0$, by differentiation of integrability.

If formality or the convergent analytic comparison is unavailable,
one must retain the full higher-order Kuranishi obstruction.
Neither a quadratic tangent obstruction nor a formal equivalence
alone justifies the homogeneous quotient and its rescalings.

\subsection{The fixed complementary gauge equation}
Use $h_0$ and its density $dV_0$ to define the Hermitian moment
\[
 \mathcal M_t=i(F^0\wedge\omega_t)/dV_0,
\]
where $F^0$ is the trace-free curvature of the unitary connection
determined by a Dolbeault operator. Let $Q_0$ be the fixed
$L^2$-orthogonal complement of the entire parallel trace-free
Hermitian algebra. The derivative in $s$ at the origin of
\[
 (t,x,s)\longmapsto
 Q_0\mathcal M_t\bigl(e^s\cdot(\bar\partial_0+a(x))\bigr)
\]
is a fixed nonzero multiple of the covariant Laplacian on $Q_0$.
All of its kernel has been removed. On fixed Sobolev spaces
$W^{k+2,2}\to W^{k,2}$, $k>2$, the implicit-function theorem gives
a unique small smooth $\sigma(t,x)\in Q_0$.
Harmonicity and the K\"ahler identities kill its first $x$ derivative
at $(0,0)$. Elliptic regularity gives, for every fixed $j$,
\begin{equation}\label{eq:complementary-correction}
 \|\sigma(t,x)\|_{C^j}\le C_j(|t|+\|x\|^2).
\end{equation}
Uniqueness gives $K$-equivariance.
Only this fixed elliptic complement is inverted.

Put $\Psi_t(x)=e^{\sigma(t,x)}\cdot(\bar\partial_0+a(x))$.
For $S\in\bigoplus_i\Herm W_i$ use the trace-free lift
\[
 \widehat S=\bigoplus_i I_{F_i}\otimes(S_i-c(S)I_{W_i}),
 \qquad c(S)=\frac{\sum_i r_i\tr S_i}{N}.
\]
Define the finite moment $m_t$ in the ordinary multiplicity dual by
\[
 \sum_i\tr(m_{t,i}(x)S_i)
   =\int_X\tr(\mathcal M_t(\Psi_t(x))\widehat S)\,dV_0.
\]
Its vanishing is equivalent to vanishing of the remaining parallel
component. This is not the rank-and-volume-weighted coefficient
of the orthogonal parallel projection. The determinant character
on $G$ is $\prod_i(\det g_i)^{r_i}$; using the ineffective common
scalar and this lift avoids an incorrect unweighted determinant
normalization.

Pullback of the affine unitary-connection symplectic form gives
a smooth closed invariant two-form $\Omega_t$ on the full ball.
In Dolbeault tangent coordinates our affine form is
$\operatorname{Im}\langle\alpha,\beta\rangle_{L^2(\omega_t)}$,
where the Hermitian pairing is linear in its first argument.
Thus the $+iF^0$ convention gives the quadratic moment $+\mu$
and $\Omega_{0,0}(iC_xS,C_xS)=\|C_xS\|^2$.
The exact moment identity is
$d\langle m_t,S\rangle=2\Omega_t(\xi_S,\cdot)$.
Since $D\Psi_0(0)$ is harmonic inclusion,
\begin{equation}\label{eq:ambient-form-estimate}
 \Omega_{0,0}=\omega_0,\qquad
 \|\Omega_{t,x}-\omega_0\|\le C(|t|+\|x\|).
\end{equation}
This proves positivity on the ordered orbit pairs used in
Section~\ref{sec:balancing}; compatibility with the original
complex structure is not needed.

At $x=0$, equivariance makes $\sigma(t,0)$ lie in
$\bigoplus_i\End F_i\otimes I_{W_i}$. It need not be scalar on $F_i$,
but preserves the factor decomposition. Chern--Weil therefore gives
\[
 \langle m_t(0),S\rangle
 =2\pi\sum_i\deg_{\omega_t}F_i\,
       \tr(S_i-c(S)I_{W_i})
 =-\sum_i\zeta_i(t)\tr S_i.
\]
The normalized moment $\rho_t=m_t-m_t(0)$ has target
$\rho_t=\zeta(t)$. Its complement has already been solved, so this
is exactly the remaining projective HYM equation.
The sign can also be checked on a single extension arrow
$A:F_i\to F_j$. Its moment is negative at $i$ and positive at $j$.
The corresponding extension $0\to F_j\to E\to F_i\to0$
can be stable only when $\mu(F_j)<\mu(E)$, precisely the sign
$\zeta_j>0$ in \eqref{eq:geometric-parameter}.
The scalar determinant equation is solved separately.
On a K3 surface $H^1(\mathcal O_X)=0$, so fixed determinant removes
no scalar degree-one deformation; its Hermitian metric does not
alter the endomorphism-valued norm.

Theorem~\ref{thm:uniform-balancing}, applied on a ball proportional
to $\|x\|$, now gives forward and reverse corrections of norm
$O(|t|+\|x\|)$ in the full multiplicity group.
They remain in $Z$ and work for every stabilizer.

\subsection{Complementary minimization and harmonic lengths}
Transport the metric to the uncorrected holomorphic family by
$H_{t,x}=h_0e^{2\sigma(t,x)}$.
Let $\mathfrak g_0$ be the full parallel complex trace-free algebra
and $Q$ its fixed $L^2$-orthogonal complement in zero-forms.
At the central point,
$\|\eta\|_{W^{1,2}}\le C\|d\eta\|_{L^2}$ on $Q$.
Smallness of $a(x)$ in $C^0$ preserves this bound with $D_x$;
nearby Hermitian and K\"ahler norms are uniformly equivalent.
Define by polarization
\begin{equation}\label{eq:complementary-minimum}
 \mathcal B_{t,x}(v,v)
 =\inf_{\eta\in QW^{1,2}}
        \|J_xv-D_x\eta\|_{H_{t,x},\omega_t}^2.
\end{equation}
The minimizing normal operator has coercive bilinear form
$\langle D_x\eta,D_x\nu\rangle$ on a fixed Hilbert space.
Its bounded inverse varies smoothly. Thus $\mathcal B$ is smooth,
positive near the origin and equivariant, and harmonic orthogonality
gives $\mathcal B_{0,0}=g_{\rm flat}$. Consequently
\begin{equation}\label{eq:complementary-metric}
 (1-C(|t|+\|x\|))g_{\rm flat}
 \le\mathcal B_{t,x}
 \le(1+C(|t|+\|x\|))g_{\rm flat}.
\end{equation}

Split every gauge zero-form into its parallel and complementary
parts. Equation~\eqref{eq:exact-gauge-derivative} gives
\begin{equation}\label{eq:full-gauge-minimum}
 \inf_\gamma\|J_xv-D_x\gamma\|^2
 =\inf_{\xi\in\mathfrak g_0}
       \mathcal B_{t,x}(v-C_x\xi,v-C_x\xi).
\end{equation}
At an integrable HYM point and a genuine stratum tangent,
Hodge theory identifies the left side with the harmonic
Kodaira--Spencer norm. A changing gauge adds an exact $D_x$ term
and conjugates isometrically, so its derivative disappears from
this minimum. We do not identify arbitrary Zariski tangent vectors
at a singular coarse point with metric tangent vectors.

For a tangent lift at a flat balanced point, differentiating the
real moment equation gives orthogonality to Hermitian orbit
directions. At a central value, Hermitian and unitary orbit
directions are orthogonal. Hence the usual quotient norm equals
\[
 \inf_{\xi\in\mathfrak g}\|v-C_x\xi\|^2
\]
on these balanced stratum tangents.
If the correction is locally $y=g(x)x$, its derivative is
$g(x)v$ plus a complex vertical vector, which disappears under
this infimum. At the point we may take the bounded $g=e^S$
from Theorem~\ref{thm:uniform-balancing}. Its squared ambient
norm distortion is between $e^{-4\|S\|_{\rm op}}$ and
$e^{4\|S\|_{\rm op}}$, and it maps complex orbit directions onto
complex orbit directions. Equations~\eqref{eq:complementary-metric}
and \eqref{eq:full-gauge-minimum} therefore yield
\eqref{eq:harmonic-comparison} without estimating $Dg$ or
inverting a small orbit eigenvalue.

A smooth local correction on each stratum exists as follows.
Fix a balanced $x_0$ with compact stabilizer $H$ and bounded
correction $S_0\in\mathfrak p_H$, and set $y_0=e^{S_0}x_0$.
Polar convexity gives $G_{x_0}=H^\CC$.
Because $S_0$ commutes with $H^\CC$, we have
$G_{y_0}=H^\CC$ and $K_{y_0}=H$.
On $V^H$ the map
\[
 (T,u)\longmapsto\rho_t(e^Te^{S_0}u)-\zeta(t),
 \qquad T\in\mathfrak p_H,
\]
takes values in that fixed space. Its derivative in $T$ at
$(0,x_0)$ has quadratic form
$2\Omega_{t,y_0}(iC_{y_0}T,C_{y_0}T)>0$ for $T\ne0$.
It is invertible at this point. The finite implicit-function
theorem gives a smooth local choice, and a compact orbit-type
slice and uniqueness identify the choices on overlaps.
At $x_0$ it is the bounded correction already used.
The local domains may shrink as orbit eigenvalues vanish; the
pointwise relative constant does not. The same argument applies
to the inverse. The holomorphic harmonic/Quot map and correction
within complex orbits preserve the complex quotient structures
on these strata.

\subsection{Full coverage and the analytic quotient}
The slice must cover all nearby actual solutions, not merely
construct a family of examples. We use the integrable HYM branch
of Delloque \cite{Del25}: the HYM subsolution condition is automatic
in \S3.2; Proposition~4.3 gives harmonic local classification,
Proposition~4.5 the complementary equation, Corollary~5.2 the
nearby morphism identification, and Proposition~5.6 the
common-parameter inverse. Lemma~7.3 and Theorem~7.4 give an
open-image homeomorphism for the integrable solution quotient.
The central solution is integrable; stable factors are simple
with vanishing cross-Hom, as required.
The possible nonintegrable failure of openness is not used.

To match its gauge normalization with ours, let $e^s$ be the
positive part of a small gauge from a harmonic representative
to an actual solution. Change the harmonic base by $e^h$ with
$h$ parallel. The new positive logarithm is
\[
 L(s,h)=\tfrac12\log(e^{-h}e^{2s}e^{-h}).
\]
Its parallel projection has derivative $-I$ in $h$ at zero.
The finite implicit-function theorem makes $L$ complementary
on a common parameter neighborhood. After unitary gauge,
uniqueness identifies it with $\sigma(t,x')$.
The reverse finite correction then provides a flat balanced
representative. Conversely every integrable zero of the corrected
moment is projectively HYM, and the determinant equation makes
it HYM. High-Sobolev classification and elliptic regularity give
the asserted smooth topology and full inverse coverage.
Nested neighborhoods are essential; no same-radius truncation
or identification of a whole noncompact quotient with a small
actual chart is claimed.

At the central value, the equivariant harmonic-to-Quot analytic
germ isomorphism descends through the analytic good quotient.
Invariant holomorphic functions pull back to invariant holomorphic
functions, so the universal property gives an isomorphism of
full reduced complex-space germs, not just their strata.
The forward and reverse moment correction identify the closed
orbits with actual HYM classes. This also applies with $H$ and
$\omega_*$ distinct under the stated hypotheses.
At a generic nonzero level the entire source chart is stable
and smooth, so the stratumwise holomorphic identification already
identifies its full complex structure.
No global comparison of nearby slope and Gieseker stability
is assumed: the latter is used at the central point for its
analytic germ.

Finally suppose a central unitary acts by $-I$ on $V$.
All ambient constructions are equivariant. The map $F$ is odd
and the coefficient matrices of $\Omega_t$ and $\mathcal B_{t,x}$
are even in $x$. Taylor expansion improves both errors to
$O(|t|+\|x\|^2)$; the same balancing and minimization proofs give
the improved rate. This completes the proof of
Theorem~\ref{thm:harmonic-comparison}.

